\chapter{Results for RQ2}
\label{cha:results-2}
\todo{Adjust labeling}

\todo{intro}

\section{Correctness tests for Algorithm~\ref{alg:find-all-k-intersection}}

To ensure that Algorithm~\ref{alg:find-all-k-intersection} works as intended and that its implementation is correct, we perform several correctness tests using known graphs. Since this algorithm is closely related to Algorithm~\ref{alg:find-k-intersection}, we use several of the same test graphs from Section~\ref{sec:1-correct}, with the expected results being the circumference of the graph and the longest-cycle intersection. For the following tests, we turn off the functionality of breaking early when during DFS on a graph $G$ a cycle is found of length $n(G)-1$ or larger.

First, we test the algorithm using a list of individual cycles on 3 to 30 vertices. We again for this test enable a setting in the program which outputs all graphs, instead of those with a longest-cycle intersection of at most 1 vertex. The program correctly identifies the circumference of every input graph and the intersection of the longest cycles as all vertices. The \texttt{graph6} string of the graph with circumference 30 in the output \texttt{CSV} file was put into the graph drawing function of the House of Graphs~\cite{data:houseofgraphs}, and this confirm that the output is the expected cycle on 30 vertices.

The program correctly identifies the Petersen graph as having circumference 9 and an empty longest-cycle intersection. Consequently, it is correctly marked as a ``solution'' in the output \texttt{CSV} file. The modified Petersen graph, obtained by inserting a vertex on each of the three edges incident to one vertex $v$, is also correctly identified as having circumference 11 and a longest-cycle intersection consisting of the single vertex $v$. It is correctly marked as a ``single-vertex intersection'' in the output \texttt{CSV} file.

Additionally, for the graphs $G_{4,2}$, $G_{5,2}$, and $G_{7,3}$, the program correctly identifies each graph as having the expected circumference 10, 13, and 25, and an empty intersection of its longest cycles. After replacing all $H_i$ in each graph with a complete graph, i.e. a graph in which every pair of distinct vertices is adjacent, the program produces the same results, as expected.

We also tested the program on the three lists of graphs from the House of Graphs~\cite{data:houseofgraphs} of 2-connected graphs on at most 32 vertices with circumference $c=10$ (533 graphs), $c=11$ (265 graphs), and $c=12$ (495 graphs). Since the second algorithm computes the circumference rather than taking $k$ as an input parameter, we expect the program to identify circumference 10 for every graph in the first list, circumference 11 for every graph in the second list, and circumference 12 for every graph in the third list. The program correctly produces these circumferences for all graphs in the respective lists. The corresponding longest-cycle intersections are also computed correctly.

Lastly, from Observation~\ref{obs:2con-n14} we know that there exists no 2-connected, non-Hamiltonian planar graph with minimum degree 4 on 13 vertices, and that there are only two such graphs on 14 vertices, each with circumference 10. We re-enable the functionality of halting enumeration early when cycles of length $n-1$ or larger are found for this test. If we run the generation pipeline mentioned in Section~\ref{sec:generators-used}, namely
\begin{verbatim}
    plantri <n> -gpm4c2x
\end{verbatim}
with \verb|<n>| set to 13, we expect to see that the number of graphs on which the DFS was halted early is equal to the number of graphs generated or checked. Similarly, with \verb|<n>| set to 14, we expect to see the same result, except for exactly 2 graphs for which we should see that these have circumference 10. This is exactly what we observe: 2800 graphs were generated and terminated early for $n=13$, and 34215 graphs were generated, of which 34213 terminated early, with 2 graphs of circumference 10 for $n=14$.

These tests give adequate confirmation that Algorithm~\ref{alg:find-all-k-intersection} was correctly designed and implemented, and functions as intended.

\section{Computational results for RQ2}

Two pipelines were used for the second research question. One of these pipelines used \texttt{plantri} to generate 2-connected planar graphs with minimum degree at least 4, while the other pipeline used a list of graphs provided to us for this thesis of non-Hamiltonian 2-connected planar graphs with minimum degree at least 4. We will once again begin with the time these pipelines took, to show that the list provided, as mentioned in Section~\ref{sec:generators-used}, greatly improved the size of graph class we can make definitive statements on.

\subsection{Timings}
\label{sec:2-timing}

Table~\ref{tab:timing-planar-m4} shows the time in seconds for the pipeline using \texttt{plantri} per vertex count, alongside the generation time by \texttt{plantri} itself and the CPU time of the program implementing Algorithm~\ref{alg:find-all-k-intersection}. Based on this growth in time for growing order of the graphs, we conclude that it is likely infeasible to generate graphs of order 20 or higher using this pipeline.

% Planar graphs
\begin{table}[ht]
\centering
\caption{Total time of the pipeline on generating 2-connected planar graphs with minimum degree at least 4.}
\label{tab:timing-planar-m4}
\begin{tabular}{rrrr}
\toprule
$n$ & Total time (s) & \texttt{plantri} CPU time (s) & Program CPU time (s) \\
\midrule
13 & 0 & 0.02 & 0.018 \\
14 & 0 & 0.24 & 0.414 \\
15 & 10 & 2.65 & 9.593 \\
16 & 309 & 31.07 & 300.739 \\
17 & 7046 & 358.74 & 6935.290 \\
\bottomrule
\end{tabular}
\end{table}

Table~\ref{tab:timing-nonhamil-up-to-19} shows the time the pipeline using the list took in seconds for graphs up to 19 vertices. The file for graphs on 20 vertices was split into 5 parts of roughly equal size. Table~\ref{tab:timing-nonhamil-20} shows the time for this pipeline on these five files for graphs on 20 vertices. 

% Non-Hamiltonian graphs up to 19 vertices
\begin{table}[ht]
\centering
\caption{Total time of the pipeline on the list of non-Hamiltonian graphs up to 19 vertices.}
\label{tab:timing-nonhamil-up-to-19}
\begin{tabular}{rrr}
\toprule
$n$ & Total time (s) & Program CPU time (s) \\
\midrule
14 & 0 & 0.002 \\
15 & 0 & 0.007 \\
16 & 0 & 0.185 \\
17 & 4 & 3.944 \\
18 & 109 & 109.156 \\
19 & 3787 & 3780.220 \\
\bottomrule
\end{tabular}
\end{table}

% Non-Hamiltonian graphs on 20 vertices
\begin{table}[ht]
\centering
\caption{Total time of the pipeline on the list of non-Hamiltonian graphs on 20 vertices.}
\label{tab:timing-nonhamil-20}
\begin{tabular}{rrr}
\toprule
Part & Total time (s) & Program CPU time (s) \\
\midrule
0 & 17007 & 17001.300 \\
1 & 17959 & 17953.800 \\
2 & 15531 & 15525.800 \\
3 & 15424 & 15417.900 \\
4 & 18029 & 18020.400 \\
\bottomrule
\end{tabular}
\end{table}

\subsection{Results and observations for RQ\ref{rq2}}

The pipeline using the list of non-Hamiltonian, 2-connected planar graphs with minimum degree at least 4 allowed computations to happen on graphs of up to 20 vertices. Per number of vertices, a table can be found in Appendix~\ref{app:rq2} of number of graphs per circumference, including per circumference the number of graphs per intersection of all longest cycles. This excludes those graphs which have a circumference equal to $n(G)-1$, though the number of graphs with this circumference is the same as the number of graphs for which the program terminated early. These numbers are also documented in Appendix~\ref{app:rq2}. This brings us to the following observation.
\begin{observation}
    \label{obs:planar-20-nothing}
    If a 2-connected planar graph on at most 20 vertices has a $k$-cycle in every vertex-deleted subgraph, then the graph contains a cycle larger than $k$.
\end{observation}

\section{New Values for RQ2}

Studying the gadget mentioned in section~\ref{sec:gadget} led to an insight in finding new negative examples to RQ\ref{rq2}. By expanding not only the inner edges between the cycles, but also expanding the edges of the cycles to a new length, graphs could be created with longest cycles that differ in length compared to the order of the graph than those already known. Similar to graphs of the form $G_{p,\ell}$, these graphs also have no vertex meeting every longest cycle. The construction of this graph depends on the integers $p$ and $\ell$, as well as on a new, ordered set of integers $D = [d_0,\ldots,d_{p-1}]$.

\begin{theorem}
  \label{theorem:d-extension}
Take any integers $p \ge 4$, $\ell \ge 2$, and integers $[d_0,\ldots,d_{p-1}] = D$ where $d_i \ge 1$ for all $i$. Let $\sigma = \sum_{i=0}^{p-1}d_i$. Then, if $\sigma \le (p-3)\ell + 2\min(D)$, there exists a 2-connected graph $G$ of order~$2\sigma + p(\ell-1)$ and circumference~$(p-1)\ell+\sigma$, whose intersection of all longest cycles is empty.
\end{theorem}
\begin{proof}
The proof takes inspiration from the proof of the similar Theorem 1 in~\cite{article:zamfirescu}. We will first construct a graph $G$, after which we will determine the types of cycles $G$ can have, alongside their length, and show that the intersection of the cycles of one type is empty. Lastly, we will compare the lengths of these cycles, and show that the empty-intersection cycles become the longest cycles per the condition.

\paragraph{Construction of $G$}
Take the graph $G_{p,\ell}$ from Section~\ref{sec:tube-graph}, with $a_i$ and $b_i$ defined as before. We use the following convention for the indices past 0 and $p-1$, to more intuitively represent the cyclical construction. For odd $p$, the continuation preserves the $a$- and $b$-cycles, giving
\[
a_{-1}=a_{p-1}, \qquad a_p=a_0,
\]
and, analogously,
\[
b_{-1}=b_{p-1}, \qquad b_p=b_0.
\]
For even $p$, the continuation crosses over to the other base cycle, giving
\[
a_{-1}=b_{p-1}, \qquad a_p=b_0,
\]
and, analogously,
\[
b_{-1}=a_{p-1}, \qquad b_p=a_0.
\]
These identifications are illustrated in Figure~\ref{fig:loop-around}.

\begin{figure}[ht]
    \centering

    %=========================================================
    % (a) p odd
    %=========================================================
    \begin{subfigure}{0.45\textwidth}
        \centering
        \begin{tikzpicture}[line cap=round,line join=round]

        %-----------------------------------------------------
        % Main vertices
        %-----------------------------------------------------
        \node[vertex] (aL) at (0,2) {};
        \node[vertex] (aR) at (2,2) {};

        \node[vertex] (bL) at (0,0) {};
        \node[vertex] (bR) at (2,0) {};

        %-----------------------------------------------------
        % Continuation edges
        %-----------------------------------------------------
        \draw[dashed,lightgray] (-0.8,2)--(aL);
        \draw[dashed,lightgray] (aR)--(2.8,2);

        \draw[dashed,lightgray] (-0.8,0)--(bL);
        \draw[dashed,lightgray] (bR)--(2.8,0);

        %-----------------------------------------------------
        % Horizontal edges
        %-----------------------------------------------------
        \draw (aL)--(aR);
        \draw (bL)--(bR);

        %-----------------------------------------------------
        % Vertical edges
        %-----------------------------------------------------
        \draw (aL)--(bL);
        \node[fill=white,inner sep=1pt,rotate=90] at (0,1) {/};

        \draw (aR)--(bR);
        \node[fill=white,inner sep=1pt,rotate=90] at (2,1) {/};

        %-----------------------------------------------------
        % Vertex labels
        %-----------------------------------------------------
        \node[above=3pt,align=center] at (0,2)
            {$a_{p-1}$\\[-2pt]\textcolor{gray}{\scriptsize $a_{-1}$}};

        \node[above=3pt,align=center] at (2,2)
            {$a_0$\\[-2pt]\textcolor{gray}{\scriptsize $a_p$}};

        \node[below=3pt,align=center] at (0,0)
            {$b_{p-1}$\\[-2pt]\textcolor{gray}{\scriptsize $b_{-1}$}};

        \node[below=3pt,align=center] at (2,0)
            {$b_0$\\[-2pt]\textcolor{gray}{\scriptsize $b_p$}};

        \end{tikzpicture}
        \caption{$p$ odd}
        \label{fig:parity-odd}
    \end{subfigure}
    \hfill
    %=========================================================
    % (b) p even
    %=========================================================
    \begin{subfigure}{0.45\textwidth}
        \centering
        \begin{tikzpicture}[line cap=round,line join=round]

        %-----------------------------------------------------
        % Main vertices
        %-----------------------------------------------------
        \node[vertex] (aL) at (0,2) {};
        \node[vertex] (aR) at (2,2) {};

        \node[vertex] (bL) at (0,0) {};
        \node[vertex] (bR) at (2,0) {};

        %-----------------------------------------------------
        % Continuation edges
        %-----------------------------------------------------
        \draw[dashed,lightgray] (-0.8,2)--(aL);
        \draw[dashed,lightgray] (aR)--(2.8,2);

        \draw[dashed,lightgray] (-0.8,0)--(bL);
        \draw[dashed,lightgray] (bR)--(2.8,0);

        %-----------------------------------------------------
        % Diagonal edges
        %-----------------------------------------------------
        \draw (aL)--(bR);
        \draw (bL)--(aR);

        %-----------------------------------------------------
        % Vertical edges
        %-----------------------------------------------------
        \draw (aL)--(bL);
        \node[fill=white,inner sep=1pt,rotate=90] at (0,1) {/};

        \draw (aR)--(bR);
        \node[fill=white,inner sep=1pt,rotate=90] at (2,1) {/};

        %-----------------------------------------------------
        % Vertex labels
        %-----------------------------------------------------
        \node[above=3pt,align=center] at (0,2)
            {$a_{p-1}$\\[-2pt]\textcolor{gray}{\scriptsize $b_{-1}$}};

        \node[above=3pt,align=center] at (2,2)
            {$a_0$\\[-2pt]\textcolor{gray}{\scriptsize $b_p$}};

        \node[below=3pt,align=center] at (0,0)
            {$b_{p-1}$\\[-2pt]\textcolor{gray}{\scriptsize $a_{-1}$}};

        \node[below=3pt,align=center] at (2,0)
            {$b_0$\\[-2pt]\textcolor{gray}{\scriptsize $a_p$}};

        \end{tikzpicture}
        \caption{$p$ even}
        \label{fig:parity-even}
    \end{subfigure}
    
    \caption{Graphical definition for indices -1 and $p$ depending on parity of $p$.}
    \label{fig:loop-around}
\end{figure}


For all $i \in \{0,\ldots,p-1\}$, insert vertices on each of the edges $(a_i,a_{i+1})$ and $(b_i,b_{i+1})$ until each is a path of length $d_i$, i.e. has a total of $d_i+1$ vertices, including the original vertices. We denote this graph as $G_{p,\ell,D}$, and denote the path on these vertices from $a_i$ to $a_{i+1}$ by $D_{a,i}$, and similarly the path from $b_i$ to $b_{i+1}$ by $D_{b,i}$. We use the same convention for the indices $p$ and $-1$, where for $p$ odd
\[
D_{a,-1}=D_{a,p-1}, \qquad D_{a,p}=D_{a,0},
\]
and
\[
D_{b,-1}=D_{b,p-1}, \qquad D_{b,p}=D_{b,0}.
\]
For even $p$, the continuation crosses over to the other base cycle, giving
\[
D_{a,-1}=D_{b,p-1}, \qquad D_{a,p}=D_{b,0},
\]
and
\[
D_{b,-1}=D_{a,p-1}, \qquad D_{b,p}=D_{a,0}.
\]
Figure~\ref{fig:g532} shows an example graph $G_{5,3,D}$ for $D=[2,2,2,2,2]$.

\begin{figure}
    \centering
    \begin{tikzpicture}
    \ExtendedTubeGraph{5}{3}{2}
    \end{tikzpicture}
    \caption{A graph $G_{p,\ell,D}$ with $p=5$, $\ell=3$,$D=[2,2,2,2,2]$.}
    \label{fig:g532}
\end{figure}

% S WITHOUT d_i:
%   S - d_i
% S WITHOUT min:
%   S - \min(D)

We also define the sets $X_i = \{a_i,b_i\}$, $A_i = \{a_i,a_{i+1}\}$, $B_i = \{b_i,b_{i+1}\}$. Since these sets are vertex-cuts, any cycle going through an internal vertex or edge of $D_{a,i}$, $D_{b,i}$, or $H_i$ (i.e. traverses these subgraphs) must also contain $A_i$, $B_i$, or $X_i$ respectively. Furthermore, the only subgraphs incident on $a_i$ are $D_{a,i-1}$, $D_{a,i}$, and $H_i$, and thus every cycle containing $a_i$ must either continue through both adjacent $D_{a,i-1}$ and $D_{a,i}$ subgraphs, or through exactly one of the aforementioned subgraphs and $H_i$. An analogous reasoning can be used for cycles going through $b_i$.
We now continue to show that there are only three types of candidate longest cycles, namely one type avoiding every $H_i$, and two types of cycles containing at least one $H_i$.

\paragraph{Cycles avoiding $H_i$}
Here we discuss the cycles avoiding every $H_i$, i.e. cycles not traversing any internal edges of any $H_i$. For $p$ odd, if we were to remove all internal vertices and edges belonging to all $H_i$, we get two disjoint cycles. As such, the longest cycles which do not traverse any $H_i$ are contained within and equal to the two subgraphs
$$
\bigcup_{i \in \{0,\dots,p-1\}}D_{a,i}\; \text{and} \;
\bigcup_{i \in \{0,\dots,p-1\}}D_{b,i}
$$
These cannot be the longest cycles of the graph, as one can create a longer cycle by picking any $i$, removing the internal vertices of $D_{a,i}$ (or $D_{b,i}$), and creating a new cycle through $H_i$, $D_{b,i}$ (or $D_{a,i}$ respectively), and $H_{i+1}$. The original cycle has length $\sigma = \sum_{i=0}^{p-1}d_i$, while this new cycle has length $\sigma+2\ell$. This new cycle must be larger than the former cycle, since by construction $\ell > 0$. As a result, any longest cycle for $p$ odd must contain at least one $H_i$.

For $p$ even, removing all internal vertices and edges belonging to all $H_i$, we get one cycle
$$
\bigcup_{i \in \{0,\dots,p-1\}}\left(D_{a,i} \cup D_{b,i}\right)
$$
with length $2\sigma$. Trivially, this is the longest cycle avoiding all $H_i$ for $p$ even.

If a cycle traverses an $H_i$, it must enter at $a_i$ from either $D_{a,i}$ or $D_{a,i-1}$, and leave at $b_i$ through either $D_{b,i}$ or $D_{b,i-1}$. We now show there are two types of cycles these options create.

\paragraph{Same-direction cycles}
A cycle can enter and leave an $H_i$ in the same direction, either through
$D_{a,i-1}$ and $D_{b,i-1}$ or through $D_{a,i}$ and $D_{b,i}$. For simplicity,
assume a cycle $C$ traverses $H_0$, $D_{a,0}$, and $D_{b,0}$. At $a_1$, there
are two possibilities: to either traverse $H_1$, closing the cycle as in Figure
\ref{fig:same-side:closed}, or to traverse $D_{a,1}$. The latter option
disallows any traversal of $H_1$, and therefore forces the choice at $b_i$ to
traverse $D_{b,1}$, as shown in Figure~\ref{fig:same-side:continued}. We then
again have the same options for traversal at $a_2$. One can see this choice of
traversing $D_{a,i}$ and $D_{b,i}$ instead of $H_i$ as extending the otherwise
closed cycle with $2d_i$ vertices.
Repeating this argument successively for $a_2$, $a_3$, \dots, $a_{p-1}$ yields
the longest same-direction cycle through $H_0$, $D_{a,0}$, and $D_{b,0}$ with
length $2(\sigma - d_{p-1})+2\ell$, and traverses all vertices

$$
\bigcup_{i \in \{0,\dots,p-2\}}V(D_{a,i})
\;\cup\;
\bigcup_{i \in \{0,\dots,p-2\}}V(D_{b,i})
\;\cup\; 
V(H_0)
\;\cup\; 
V(H_{p-1})
$$

The reasoning for $D_{a,-1}$ and $D_{b,-1}$, and for other $H_i$'s, is
analogous. The largest of these cycles is all cycles which avoid a subgraph pair $D_{a,i}$ and $D_{b,i}$ of minimal length, i.e. where $d_i = \min(D)$, with a cycle length of $2(\sigma - \min(D))+2\ell$. Note that all vertices $a_i$ and $b_i$ meet all these cycles.

\begin{figure}
    \centering

    %=========================================================
    % Version (1)
    %=========================================================
    \begin{subfigure}{0.45\textwidth}
    \centering
    \resizebox{\linewidth}{!}{
    \begin{tikzpicture}[line cap=round,line join=round]

    %---------------------------------------------------------
    % Main vertices
    %---------------------------------------------------------
    \node[vertex,draw=red,label=above:$a_0$] (a0) at (0,2) {};
    \node[vertex,draw=red,label=above:$a_1$] (a1) at (2,2) {};
    \node[vertex,draw=red,label=above:$a_2$] (a2) at (4,2) {};

    \node[vertex,draw=red,label=below:$b_0$] (b0) at (0,0) {};
    \node[vertex,draw=red,label=below:$b_1$] (b1) at (2,0) {};
    \node[vertex,draw=red!40,label=below:$b_2$] (b2) at (4,0) {};

    %---------------------------------------------------------
    % Continuation edges
    %---------------------------------------------------------
    \draw[dashed,lightgray] (-0.8,2)--(a0);
    \draw[dashed,lightgray] (a2)--(4.8,2);

    \draw[dashed,lightgray] (-0.8,0)--(b0);
    \draw[dashed,lightgray] (b2)--(4.8,0);

    %---------------------------------------------------------
    % Horizontal edges
    %---------------------------------------------------------
    % D_{a,0}
    \draw[red,very thick] (a0)--(a1);
    \node[fill=white,inner sep=1pt,text=red] at (1,2) {//};

    % D_{a,1}
    \draw[red,very thick] (a1)--(a2);
    \node[fill=white,inner sep=1pt,text=red] at (3,2) {///};

    % D_{b,0}
    \draw[red,very thick] (b0)--(b1);
    \node[fill=white,inner sep=1pt,text=red] at (1,0) {//};

    % D_{b,1}
    \draw[red!40,very thick] (b1)--(b2);
    \node[fill=white,inner sep=1pt,text=red!40] at (3,0) {///};

    %---------------------------------------------------------
    % Vertical edges
    %---------------------------------------------------------
    % H_0
    \draw[red,very thick] (a0)--(b0);
    \node[fill=white,inner sep=1pt,rotate=90,text=red] at (0,1) {/};

    % H_1
    \draw (a1)--(b1);
    \node[fill=white,inner sep=1pt,rotate=90] at (2,1) {/};

    % H_2
    \draw (a2)--(b2);
    \node[fill=white,inner sep=1pt,rotate=90] at (4,1) {/};

    \end{tikzpicture}
    }
    \caption{}
    \label{fig:same-side:continued}
    \end{subfigure}
    \hfill
    %=========================================================
    % Version (2)
    %=========================================================
    \begin{subfigure}{0.45\textwidth}
    \centering
    \resizebox{\linewidth}{!}{
    \begin{tikzpicture}[line cap=round,line join=round]

    %---------------------------------------------------------
    % Main vertices
    %---------------------------------------------------------
    \node[vertex,draw=red,label=above:$a_0$] (a0) at (0,2) {};
    \node[vertex,draw=red,label=above:$a_1$] (a1) at (2,2) {};
    \node[vertex,label=above:$a_2$] (a2) at (4,2) {};

    \node[vertex,draw=red,label=below:$b_0$] (b0) at (0,0) {};
    \node[vertex,draw=red,label=below:$b_1$] (b1) at (2,0) {};
    \node[vertex,label=below:$b_2$] (b2) at (4,0) {};

    %---------------------------------------------------------
    % Continuation edges
    %---------------------------------------------------------
    \draw[dashed,lightgray] (-0.8,2)--(a0);
    \draw[dashed,lightgray] (a2)--(4.8,2);

    \draw[dashed,lightgray] (-0.8,0)--(b0);
    \draw[dashed,lightgray] (b2)--(4.8,0);

    %---------------------------------------------------------
    % Horizontal edges
    %---------------------------------------------------------
    % D_{a,0}
    \draw[red,very thick] (a0)--(a1);
    \node[fill=white,inner sep=1pt,text=red] at (1,2) {//};

    % D_{a,1}
    \draw (a1)--(a2);
    \node[fill=white,inner sep=1pt] at (3,2) {///};

    % D_{b,0}
    \draw[red,very thick] (b0)--(b1);
    \node[fill=white,inner sep=1pt,text=red] at (1,0) {//};

    % D_{b,1}
    \draw (b1)--(b2);
    \node[fill=white,inner sep=1pt] at (3,0) {///};

    %---------------------------------------------------------
    % Vertical edges
    %---------------------------------------------------------
    % H_0
    \draw[red,very thick] (a0)--(b0);
    \node[fill=white,inner sep=1pt,rotate=90,text=red] at (0,1) {/};

    % H_1
    \draw[red,very thick] (a1)--(b1);
    \node[fill=white,inner sep=1pt,rotate=90,text=red] at (2,1) {/};

    % H_2
    \draw (a2)--(b2);
    \node[fill=white,inner sep=1pt,rotate=90] at (4,1) {/};

    \end{tikzpicture}
    }
    \caption{}
    \label{fig:same-side:closed}
    \end{subfigure}

    \caption{A graph segment of $G_{p,\ell,D}$ with the two possibilities at $a_1$ for cycle through $H_0$, $D_{a,0}$, and $D_{b,0}$.}
    \label{fig:same-side}
\end{figure}

\paragraph{Opposite-direction cycles}
The only other option for a cycle traversing $H_i$ is to enter and leave
$H_i$ in opposite directions, either through $D_{a,i-1}$ and $D_{b,i}$ or
through $D_{a,i}$ and $D_{b,i-1}$.

Say a cycle traversed $H_0$, $D_{a,0}$, and $D_{b,-1}$. There are again two possibilities at $a_1$: to traverse $D_{a,1}$, or to traverse $H_1$. In the latter option, the traversal of $D_{b,1}$ is forced. This is because the only other option, to traverse $D_{b,0}$, would result in the same-direction case from before, which results in a cycle where an $H_i$ being exited in opposite directions is not possible. 

From then on, this option repeats: at any endpoint $b_i$, either $D_{b,i}$ is traversed, or $H_i$ and $D_{a,i}$ are traversed, and analogous for endpoints $a_i$, up until $i=p-1$. There, $D_{b,-1}$ must be traversed to complete the cycle, either from $a_{-1}$ through $H_{p-1}$, or directly from $b_{-1}$.\footnote{Remember that $a_{-1} = b_{p-1}$ if $p$ is even} In total, every $H_i$ traversed this way adds $d_i+\ell$ vertices to the cycle (including $H_0$, where $b_0$ is counted from the traversal of $D_{b,-1}$), while every $H_i$ avoided adds only $d_i$ vertices.

By construction, a cycle in $G$ cannot traverse all $H_i$. Since $X_i$ is a vertex-cut in $G$, the subgraph $H_i$ can only be traversed by entering through $a_i$ and exiting through $b_i$. For a cycle to traverse all $H_i$, it must visit an even number of $H_i$'s if $p$ is odd to close the cycle on the same base cycle it started, or an odd number of $H_i's$ if $p$ is even, a contradiction. Thus, at least one $H_i$ must be avoided, which means the largest opposite-side cycle must traverse (1) a maximal amount of $H_i$'s, and (2) at most all but one $H_i$. We can construct a cycle $C$ avoiding exactly one $H_i$, with length $(p-1)\ell + \sigma$. Put
\[
C=
\left(
\bigcup_{i\neq 0} H_i
\right)
\cup
\left(
\bigcup_{\substack{0\leq j\leq p-1\\[2pt] j\ \mathrm{even}}} D_{a,j}
\right)
\cup
\left(
\bigcup_{\substack{0\leq k\leq p-1\\[2pt] k\ \mathrm{odd}}} D_{b,k}
\right)
\]
where $C$ omits all vertices $V(H_0)\setminus\{b_0\}$ and $V(D_{b,0})\setminus\{b_1\}$. For all other vertices $v \in V(G)$, a cycle of length $(p-1)\ell + \sigma$ can be constructed in the same way. Figure~\ref{fig:opposite-side:g532} shows this cycle for a graph $G_{5,3,D}$, and Figure~\ref{fig:opposite-side:g632} shows this cycle for a graph $G_{6,3,D}$ with $D=[2,2,2,2,2]$.

\begin{figure}
  \centering
  \begin{tikzpicture}
    \ExtendedTubeGraph{5}{3}{2}

    \foreach \v in {
    A1,DA1-1,A2,
    H2-1,H2-2,B2,
    DB2-1,B3,
    H3-2,H3-1,A3,
    DA3-1,A4,
    H4-1,H4-2,B4,
    DB4-1,B5,
    H5-2,H5-1,A5,
    DA5-1}{
    \node[vertex,draw=red,fill=white] at (\v) {};
    }
    % Highlighted cycle
    \draw[red, very thick]
      (A1) --
      (DA1-1) --
      (A2) --
      (H2-1) --
      (H2-2) --
      (B2) --
      (DB2-1) --
      (B3) --
      (H3-2) --
      (H3-1) --
      (A3) --
      (DA3-1) --
      (A4) --
      (H4-1) --
      (H4-2) --
      (B4) --
      (DB4-1) --
      (B5) --
      (H5-2) --
      (H5-1) --
      (A5) --
      (DA5-1) --
      (A1);

  \end{tikzpicture}
  \caption{A longest opposite-direction cycle marked red in a graph $G_{5,3,D}$ with $D=[2,2,2,2,2]$.}
  \label{fig:opposite-side:g532}
\end{figure}

\begin{figure}
  \centering
  \begin{tikzpicture}
    \ExtendedTubeGraph{6}{3}{2}

    \foreach \v in {
    A1,DA1-1,A2,
    H2-1,H2-2,B2,
    DB2-1,B3,
    H3-2,H3-1,A3,
    DA3-1,A4,
    H4-1,H4-2,B4,
    DB4-1,B5,
    H5-2,H5-1,A5,
    DA5-1,A6,
    H6-1,H6-2,
    B6,DB6-1}{
    \node[vertex,draw=red,fill=white] at (\v) {};
    }
    % Highlighted cycle
    \draw[red, very thick]
      (A1) --
      (DA1-1) --
      (A2) --
      (H2-1) --
      (H2-2) --
      (B2) --
      (DB2-1) --
      (B3) --
      (H3-2) --
      (H3-1) --
      (A3) --
      (DA3-1) --
      (A4) --
      (H4-1) --
      (H4-2) --
      (B4) --
      (DB4-1) --
      (B5) --
      (H5-2) --
      (H5-1) --
      (A5) --
      (DA5-1) --
      (A6) --
      (H6-1) --
      (H6-2) --
      (B6) --
      (DB6-1) --
      (A1);

  \end{tikzpicture}
  \caption{A longest opposite-direction cycle marked red in a graph $G_{6,3,D}$ with $D=[2,2,2,2,2]$.}
  \label{fig:opposite-side:g632}
\end{figure}

\paragraph{Length comparison}
We have determined there are three types of cycles with the potential of being the longest cycles, the length of which are $2\sigma$ for $H_i$-avoiding cycles when $p$ is even, $2(\sigma - \min(D))+2\ell$ for same-direction cycles and $(p-1)\ell + \sigma$ for opposite-direction cycles. Only in the last case can we construct a cycle for every vertex avoiding that vertex. The opposite-side cycle is longer than the same-side cycle if
\[
\begin{aligned}
2(\sigma-\min(D))+2\ell
    &\le (p-1)\ell + \sigma\\
\iff\qquad
\sigma
    &\le (p-3)\ell + 2\min(D)
\end{aligned}
\]
while the opposite-side cycle is longer than the $H_i$ avoiding cycle if
\[
\begin{aligned}
2\sigma
    &\le (p-1)\ell + \sigma\\
\iff\qquad
\sigma
    &\le (p-1)\ell
\end{aligned}
\]

We can show that the condition that $\sigma \le (p-3)\ell + 2\min(D)$ implies the condition $\sigma \le (p-1)\ell$.

First assume that $\min(D) \le \ell$. Then $(p-3)\ell + 2\min(D) \le (p-3)\ell + 2\ell = (p-1)\ell$, thus if the first condition is true, the second must be as well. 

Now assume instead that $\min(D) > \ell$. Then, $(p-3)\ell + 2\min(D) > (p-1)\ell$. However, we also have that $\sigma \ge p \min(D) = (p-3)\min(D) + 3\min(D)$. Since for all $d_i$ we have that $d_i\ge1$, we know that $\sigma > (p-3)\min(D) + 2\min(D)$, which after applying the assumptions gives us $\sigma > (p-3)\ell + 2\min(D)$ and, since $\min(D) > \ell$, we have $\sigma > (p-1)\ell$. This means that in this case, neither assumption can be true.

\paragraph{Conclusion}
As a result, there exists a 2-connected graph $G$ of order $2\sigma + p(\ell-1)$ with circumference $(p-1)\ell + \sigma$ for which no vertex meets every longest cycle if $\sigma \le (p-3)\ell + 2\min(D)$.

\end{proof}

Theorem~\ref{theorem:d-extension} provides us with new graphs that we can expand on in the same way as done in~\cite{article:zamfirescu} to find new negative values for RQ\ref{rq2}. First, observe that all graphs $G_{p,\ell,D}$ are planar if $p$ is odd. To prove the existence of new negative values, we must show we can add edges to some of these graphs to ensure minimum degree 4.

\begin{theorem}
  \label{theorem:new-vals-rq2}
  Take integers $p\ge5$ odd, $\ell\ge5$, and integers $[d_0,\dots,d_{p-1}]=D$ where $d_i = 1$ or $d_i\ge5$ with $\sigma=\sum_{i=0}^{p-1}d_i$. If $\sigma \le (p-3)\ell + 2\min(D)$, there exists a 2-connected, planar graph of minimum degree~4 with order~$2\sigma + p(\ell-1)$ and circumference~$k=(p-1)\ell + \sigma$, in which every vertex-deleted subgraph contains a $k$-cycle.
\end{theorem}
\begin{proof}
The proof takes inspiration from the proof of the similar Proposition 1 in~\cite{article:zamfirescu}.
Let $G$ be a graph $G_{p,\ell,D}$. Observe that in each $D_{a,i}$, $D_{b,i}$, and $H_i$, we may add edges as desired within each respective subgraph (but without adding edges between internal nodes of different subgraphs) without altering the fact that the intersection of all longest cycles remains empty. We may ignore the subgraphs $D_{a,i}$, $D_{b,i}$ for which $d_i = 1$, as these have no internal vertices, and the following construction will ensure the vertices $a_i, a_{i+1}$ and $b_i, b{i+1}$ will have a degree at least 4. We denote the modified subgraphs with additional edges as $D_{a,i}'$, $D_{b,i}'$, respectively $H_i'$.

Trivially, if we choose $D_{a,i}'$ and $D_{b,i}'$ to each be the path $P_{d_i}^2$, i.e. a path of length $d_i$ with an additional edge between all vertices at distance 2 as in Figure~\ref{fig:path-2}, and similarly choose $H_i'$ to be the path $P_\ell^2$, then $G$ remains planar. Observe that, for each of these subgraphs, two internal vertices, say $w$ and $w'$, have degree 3, while all other internal vertices have degree 4. These vertices $w$ and $w'$ are adjacent to an endpoint of the (modified) paths, say $x$ and $y$ respectively. Adding the edges $w'x$ and $wy$ when $d_i$ (respectively $\ell$) is odd as in Figure~\ref{fig:degree-4-opposite-side}, and adding the edge $ww'$ when $d_i$ (respectively $\ell$) is even as in Figure~\ref{fig:degree-4-same-side}, provides a graph $G$ with minimum degree 4 when $\ell \ge 5$ and, for all $i$, $d_i = 1$ or $d_i \ge 5$.

We conclude with the statement that for the resulting graph $G$, the intersection of all longest cycles is empty. Theorem~\ref{theorem:intersection-vertex-deletion} has shown that this empty intersection is equivalent to claiming the graph $G$ has a $k$-cycle in every vertex-deleted subgraph, and no larger cycle than $k$ exists in $G$. As such, we have shown that $G$ is a planar graph with minimum degree 4, with a $k$-cycle in every vertex-deleted subgraph, with no cycle larger than $k$.

\begin{figure}[ht]
  \centering

  % First graph
  \begin{tikzpicture}[>=stealth]
    % Nodes
    \node[vertex, label=left:$x$] (v1) at (0,0) {};
    \node[vertex] (v2) at (1,1) {};
    \node[vertex] (v3) at (2,0) {};
    \node[vertex] (v4) at (3,1) {};
    \node[vertex] (v5) at (4,0) {};
    \node[vertex, label=right:$y$] (v6) at (5,1) {};

    % Path edges
    \draw (v1) -- (v2) -- (v3) -- (v4) -- (v5) -- (v6);

    % Edges between vertices at distance 2
    \draw (v1) -- (v3);
    \draw (v2) -- (v4);
    \draw (v3) -- (v5);
    \draw (v4) -- (v6);
  \end{tikzpicture}
  \caption{A path $P_5^2$ with endpoints $x$ and $y$.}
  \label{fig:path-2}
\end{figure}

\begin{figure}
  \centering
  \begin{tikzpicture}[>=stealth]
    % Nodes
    \node[vertex, label=left:$x$] (v1) at (0,0.5) {};
    \node[vertex] (v2) at (1,1) {};
    \node[vertex] (v3) at (2,0) {};
    \node[vertex] (v4) at (3,1) {};
    \node[vertex] (v5) at (4,0) {};
    \node[vertex, label=right:$y$] (v6) at (5,0.5) {};

    % Path edges
    \draw (v1) -- (v2) -- (v3) -- (v4) -- (v5) -- (v6);

    % Edges between vertices at distance 2
    \draw (v1) -- (v3);
    \draw (v2) -- (v4);
    \draw (v3) -- (v5);
    \draw (v4) -- (v6);

    % Additional endpoint edges
    \draw[lightgray] (v1) -- (-0.5,-0.5);
    \draw[lightgray] (v1) -- (-0.5,1.5);
    \draw[lightgray] (v6) -- (5.5,-0.5);
    \draw[lightgray] (v6) -- (5.5,1.5);

    % Additional red edges
    \draw[red] (v2) to[bend left=25] (v6);
    \draw[red] (v5) to[bend left=25] (v1);
  \end{tikzpicture}
  \caption{The additional edges (marked in red) placed on either $D_{a,i}'$, $D_{b,i}'$, or $H_i'$ when $d$ or $l$ $=5$ odd.}
  \label{fig:degree-4-opposite-side}
\end{figure}


\begin{figure}
  \centering
  \begin{tikzpicture}[>=stealth]
    % Nodes
    \node[vertex, label=left:$x$] (v1) at (0,0.5) {};
    \node[vertex] (v2) at (1,1) {};
    \node[vertex] (v3) at (2,0) {};
    \node[vertex] (v4) at (3,1) {};
    \node[vertex] (v5) at (4,0) {};
    \node[vertex] (v6) at (5,1) {};
    \node[vertex, label=right:$y$] (v7) at (6,0.5) {};

    % Path edges
    \draw (v1) -- (v2) -- (v3) -- (v4) -- (v5) -- (v6) -- (v7);

    % Edges between vertices at distance 2
    \draw (v1) -- (v3);
    \draw (v2) -- (v4);
    \draw (v3) -- (v5);
    \draw (v4) -- (v6);
    \draw (v5) -- (v7);

    % Additional endpoint edges
    \draw[lightgray] (v1) -- (-0.5,-0.5);
    \draw[lightgray] (v1) -- (-0.5,1.5);
    \draw[lightgray] (v7) -- (6.5,-0.5);
    \draw[lightgray] (v7) -- (6.5,1.5);

    % Additional red edge
    \draw[red] (v2) to[bend left=25] (v6);
  \end{tikzpicture}
  \caption{The additional edges (marked in red) placed on either $D_{a,i}'$, $D_{b,i}'$, or $H_i'$ when $d$ or $l$ $=6$ even.}
  \label{fig:degree-4-same-side}
\end{figure}

\end{proof}

\section{Conclusion}
\todo{conclusion}
The final section of the chapter gives an overview of the important results
of this chapter. This implies that the introductory chapter and the
concluding chapter don't need a conclusion.
